% 鸟神星（综述）
% license CCBYSA3
% type Wiki

本文根据 CC-BY-SA 协议转载翻译自维基百科\href{https://en.wikipedia.org/wiki/Makemake}{相关文章}。

\begin{figure}[ht]
\centering
\includegraphics[width=6cm]{./figures/abca2436367ca8e9.png}
\caption{2015 年 4 月，哈勃空间望远镜拍摄的低分辨率阋神星及其无名卫星 S/2015 (136472) 1 图像} \label{fig_Makema_1}
\end{figure}
玛克马克$^{[g]}$（小行星编号：136472 Makemake）是一颗位于柯伊伯带的矮行星，柯伊伯带是海王星轨道之外由冰质天体组成的圆盘。它是第四大海王星外天体，也是经典柯伊伯带中体积最大的成员$^{[h]}$，其直径约为冥王星的 60\%。它于 2005 年 3 月 31 日由美国天文学家迈克尔·E·“迈克”·布朗（Michael E. "Mike" Brown）、查德·特鲁希略（Chad Trujillo）和大卫·拉比诺维茨（David Rabinowitz）在帕洛马天文台发现。作为该团队发现的最大天体之一，玛克马克的发现促成了 2006 年冥王星被重新归类为矮行星。

玛克马克在表面特征方面与冥王星类似：其表面高度反射，大部分覆盖着冻结的甲烷，并因凝聚物（tholins）而被染成红褐色$^{[i]}$。玛克马克目前已知有一颗卫星，但尚未命名。该卫星的轨道表明玛克马克自转具有很高的轴倾角，这意味着它会经历极端季节。玛克马克显示出地球化学活动与低温火山活动的证据，这使科学家怀疑其可能拥有一个地下液态水海洋。在玛克马克上发现了气态甲烷，但尚不清楚它是存在于大气中，还是来自暂时性的气体逸出。

由于尚未有太空探测器近距离访问玛克马克，因此并不存在关于其表面的高分辨率影像。玛克马克距离地球非常遥远，即使通过望远镜观测，它看起来也只是一颗类似星点的光斑。科学家表达了希望派遣太空探测器访问玛克马克的愿望，因为它具有地质活动及潜在的地下海洋。
\subsection{历史}
\subsubsection{发现}
\begin{figure}[ht]
\centering
\includegraphics[width=8cm]{./figures/7d2c811983023908.png}
\caption{玛克马克是在帕洛马天文台使用 1.22 米塞缪尔·奥辛望远镜（如图）拍摄的图像中被发现的} \label{fig_Makema_2}
\end{figure}
玛克马克于 2005 年由美国天文学家迈克尔·E·“迈克”·布朗（Michael E. "Mike" Brown）、查德·特鲁希略（Chad Trujillo）和大卫·拉比诺维茨（David Rabinowitz）组成的团队发现，他们当时正在寻找海王星轨道之外的大型天体$^{[23]:200}$。该团队对海王星外天体的搜索始于 2001 年$^{[23]:187-188}$，具体做法是使用安装在加利福尼亚州帕洛马天文台塞缪尔·奥辛 1.22 米（48 英寸）望远镜上的 CCD（电荷耦合器件）相机$^{[j]}$，定期对夜空进行成像$^{[25][23]:190}$。玛克马克的发现图像由该望远镜于 2005 年 3 月 31 日拍摄$^{[1][26]}$，但直到 2005 年 4 月 3 日，迈克·布朗在检查这些图像时才发现该天体，并识别出它具有异常高的亮度$^{[27]:132}$。

在发现玛克马克前几个月，布朗及其团队还发现了两个异常巨大的海王星外天体，闇神星（Haumea）与厄里斯（Eris），它们被认为至少与当时的第九行星冥王星体积相当$^{[28]}$。当团队正计划对这两个天体开展进一步观测时，原本计划将玛克马克的公布推迟至 2005 年 10 月厄里斯公布之后$^{[23]:202[27]:133-134}$。然而这一计划被打乱，因为 2005 年 7 月 27 日，由西班牙内华达山天文台 José Luis Ortiz Moreno 领导的团队宣布他们自己的闇神星发现结果$^{[27]:145[23]:207}$。布朗意识到，他所在团队包含闇神星、厄里斯与玛克马克位置的观测日志无意中公开，并被 Ortiz 所在机构的一台计算机访问过$^{[27]:154-155[23]:211}$。担心厄里斯与玛克马克的发现也会被抢先公布，布朗于 2005 年 7 月 29 日联系小行星中心（MPC）的 Brian G. Marsden 宣布发现$^{[29][27]:156}$。MPC 在加利福尼亚当地时间中午在其网站上发布了厄里斯与玛克马克的发现公告，随后当天晚上由中央天文电报局发布消息$^{[23]:210[30][26]}$。这些与冥王星大小相当的天体的公布，引发了关于“什么是行星”的广泛争论$^{[25]}$，这促使国际天文学联合会（IAU）在 2006 年 8 月制定新的行星定义，将冥王星重新归类为矮行星$^{[31][32]}$。
\subsubsection{名称与符号}
这颗矮行星以玛克马克命名，玛克马克是复活节岛本土拉帕努伊人神话中人类创造者与丰饶之神$^{[6]}$。其小行星目录编号为 136472，该编号由小行星中心在 2005 年 9 月 7 日授予，当时其轨道已被充分确定$^{[33][34]}$。在命名之前，玛克马克的临时编号为 2005 FY9，当其发现被公布时由小行星中心赋予$^{[30][6]}$。玛克马克之前还被其发现团队昵称为“Easterbunny”$^{[k]}$，这是引用其发现时间（复活节之后不久），并曾被布朗用于发现记录的自动代码名称“ K05331A ”$^{[23][7]}$。

在其个人文字与访谈中，布朗回忆称，确定玛克马克的命名非常困难，因为当时已知的天体特征与神话无明显联系$^{[7][27]:246[35][36]}$。为了保留该天体与复活节的联系，布朗曾考虑用盎格鲁-撒克逊女神 Ēostre，或阿尼什纳贝族的诡计兔 Manabozho 命名，但发现这两个名称都无法使用$^{[l]}$。最终，布朗与其团队选择了“Makemake”这一名称，满足了天体与复活节的联系，以及 IAU 关于以创世神命名经典柯伊伯带天体的规定$^{[32][37]}$。玛克马克的名称于 2008 年 7 月获 IAU 批准并公布$^{[m]}$。

玛克马克的符号 ⟨🝼⟩ 于 2022 年 1 月作为 U+1F77C 被纳入 Unicode$^{[38]}$。IAU 不鼓励在科学出版物中使用行星符号$^{[39]}$，因此该符号主要由占星学者使用$^{[40]}$。不过 NASA 曾在 2015 年发布的一份信息图中使用过该符号一次$^{[41][40]:4}$。玛克马克符号由 Denis Moskowitz 与 John T. Whelan 设计，它代表玛克马克面部的传统岩刻画，并经过设计使其类似字母“M”$^{[42]}$。其他占星师也设计并使用自己的玛克马克符号，例如 ⟨⟩$^{[40]:5}$。
\subsection{轨道与分类}
\begin{figure}[ht]
\centering
\includegraphics[width=8cm]{./figures/bc08e081cf518854.png}
\caption{示意图显示了玛克马克绕太阳运行的倾斜轨道（灰色），并标示了外行星的位置。沿着玛克马克轨道路径的垂直灰线标出其位于黄道平面之上和之下的位置。} \label{fig_Makema_3}
\end{figure}
玛克马克在海王星之外绕太阳运行，其平均距离为 45.5 个天文单位（AU；68.1 亿公里或 42.3 亿英里）$^{[b]}$。它完成一次公转需要 307 年$^{[b]}$。玛克马克的轨道偏心率为 0.16，遵循一条中等椭圆轨道$^{[44]:212-213}$，距离太阳最近时为 38.2 AU（近日点），最远时为 52.8 AU（远日点）$^{[b]}$。玛克马克相对于黄道的轨道倾角较大，为 29°$^{[9][44]:212}$。

玛克马克目前接近远日点，即其轨道上距离最远的位置$^{[16]:9}$。截至 2025 年 11 月，它距离太阳为 52.7 AU$^{[45][21]}$，并将于 2033 年 5 月抵达远日点$^{[46]}$。玛克马克目前位于黄道平面之上很远的位置$^{[47][44]:212}$，并将在远日点保持这一状态，届时其黄道纬度将达到 25.9°$^{[46]}$。玛克马克将在 2103 年穿过黄道平面$^{[48]}$，并将在 2186 年到达近日点，位于黄道平面以下 26°$^{[11]}$。N 体模拟显示，玛克马克的轨道在数十亿年尺度上是稳定的，在太阳系余下寿命中不太可能发生显著变化$^{[49]:6-7}$。

玛克马克与海王星之外许多小型冰质天体共享其轨道特征，这些天体共同构成了柯伊伯带这一区域。玛克马克特别属于“动力学高温”的经典柯伊伯带天体群$^{[n][44]:212}$，其轨道特征具有较高的倾角（i > 5°）、较低的偏心率（e < 0.2），且与海王星不存在轨道共振$^{[50]:21[51]:2}$。玛克马克是经典柯伊伯带中体积最大的成员$^{[44]:212}$，尽管它仅占该带总质量的一小部分$^{[52]:8[o]}$。动力学高温的经典柯伊伯带天体被认为在太阳系早期历史中被海王星引力散射$^{[50]:22}$，因此天文学家也将玛克马克称为“散射”天体$^{[53][54]:L98[55]:284}$。

科学界的共识是玛克马克是一颗矮行星：即它质量足够大，其自身引力可以使其形成接近球形的形状，但质量又不足以清除轨道上的其他天体，这一点从其位于柯伊伯带即可证明$^{[37][32]}$。玛克马克是 IAU 在新程序下命名的首个预期成为矮行星的天体，也是 2006 年确立该分类以来继原有的谷神星、冥王星与厄里斯之后，第四个被宣布为矮行星的天体$^{[m]}$。玛克马克更具体地说是一颗冥族矮行星（plutoid），即绕海王星之外运行的矮行星子类$^{[6][35]}$。
\subsection{大小、形状与质量}
\begin{figure}[ht]
\centering
\includegraphics[width=8cm]{./figures/c883c7e09b0f9dea.png}
\caption{比较直径大于 700 公里（430 英里）的各大海王星外天体的尺寸、反照率与颜色。玛克马克位于第一行，从右数第二个。深色弧线表示天体尺寸的不确定范围。} \label{fig_Makema_4}
\end{figure}
玛克马克是一个近乎球形的天体，其平均直径约为 1,430 公里（890 英里）$^{[12]:2}$，约为冥王星直径的 60\%（3/5）$^{[p][57]}$，或地球直径的 11\%（1/9）$^{[31]}$。这使得玛克马克成为太阳系中已知的第四大矮行星及海王星外天体，仅次于冥王星、厄里斯和闇神星$^{[58]}$。2011 年对一次恒星掩食的观测显示，玛克马克在极地略呈扁平状，极直径上限约为 1,420 公里（880 英里）$^{[q]}$，赤道直径约为 1,434 公里（891 英里）$^{[12]}$。这些尺寸与一种称为麦克劳林扁球（Maclaurin spheroid）的扁平球形结构相一致，这种形状出现在天体处于流体静力平衡状态（即天体自身引力足以将其压成球形）并因自转而发生形变时$^{[12]:2[17]:5[19]:12}$。

玛克马克的质量估计在约 (2.5\times10^{21}) 至 (2.9\times10^{21}) 千克之间，该数值由其卫星的轨道周期与距离确定$^{[15]:3}$。这使得玛克马克成为太阳系中已知质量第四大的矮行星及海王星外天体，仅次于厄里斯、冥王星和闇神星$^{[59]}$。与太阳系其他天体相比，玛克马克的质量约为地球月球质量的 3.7\%（或地球质量的 0.045\%）$^{[r]}$，约为冥王星质量的 20\%$^{[s]}$。根据玛克马克的质量和平均直径，其平均表面重力约为 0.35 m/s(^2)$^{[e]}$（约为地球重力的 3.6\%）$^{[t]}$，其表面逃逸速度约为 0.71 km/s$^{[f][16]:8[ ]}$。

\subsection{自转}
玛克马克的自转周期尚不确定，截至 2025 年的测量显示，该周期为 11.4 或 22.8 小时（0.48 或 0.95 日）$^{[19]:2,7}$。这些自转周期的测量通过监测玛克马克亮度随时间的变化进行，并绘制其光变曲线$^{[17][19]:2}$。玛克马克的亮度变化非常微小（0.03 星等），推测是由于其表面反照率变化较小，这使得望远镜很难测量玛克马克的光变曲线与自转周期$^{[17]:1,6}$。例如，2019 年之前的研究提出了可能的自转周期 7.77、11.24、11.5 与 22.48 小时$^{[17]:1}$。截至 2025 年的测量尚不清楚玛克马克的亮度在一次自转中是出现一次还是两次峰值，因此不确定其自转周期究竟是 11.4 小时，还是其两倍即 22.8 小时$^{[19]:2}$。

玛克马克的自转轴倾角尚未被测量，不过可以合理推测其自转轴与其卫星轨道的极点对齐$^{[8]:4[61]:16}$。在这种情况下，玛克马克的自转轴相对于绕太阳轨道的倾角介于 46° 至 78° 之间（或相对于黄道平面介于 63°–87° 之间），在其卫星被发现时，其赤道朝向太阳与地球（接近春分）$^{[8]:3-4}$。如此高的轴倾角加上其偏心轨道，可导致玛克马克表面温度与地形出现显著季节变化，类似于冥王星$^{[8]:4-5[61]:16}$。玛克马克的卫星被预测将在 2009–2013 年或 2023–2027 年期间掩食玛克马克，因此如果其自转与其卫星轨道一致，玛克马克可能已经在这两个年份范围之一经历过春分$^{[15]:1}$。
\subsection{地质}
\subsubsection{表面}
\begin{figure}[ht]
\centering
\includegraphics[width=10cm]{./figures/a8a10a028dc3e3f6.png}
\caption{詹姆斯·韦布空间望远镜测得的玛克马克近红外光谱。甲烷（$CH_4$）在 1–2 微米范围内的吸收特征在玛克马克的光谱中极为显著，说明其表面甲烷含量非常丰富。玛克马克上检测到的其他化合物包括乙烷（$C_{2}H_6$）、乙炔（$C_2H_2$）、氘代甲烷（$CH_3D$），以及可能存在的乙烯（$C_2H_4$）$^{[16]:2}$。} \label{fig_Makema_5}
\end{figure}
\begin{figure}[ht]
\centering
\includegraphics[width=8cm]{./figures/6453132cf5651cb1.png}
\caption{玛克马克的艺术示意图，呈现其均匀的浅棕色表面以及缺乏显著大气的特征} \label{fig_Makema_7}
\end{figure}